\section{Specific Objectives}
\label{sec:objectives}

\begin{enumerate}

\item To analyze the structure of the \href{https://epic-kitchens.github.io/}{\textit{EPIC-KITCHENS-100}} dataset and the relevance matrices used in the multimodal retrieval challenge, identifying the organization of video segments, textual descriptions, and their semantic relationships, producing a structured characterization covering segment count, label distribution, and relevance matrix statistics, that will inform preprocessing and training decisions \textbf{(Weeks 1--2)}.

\item To design a dual-encoder deep learning architecture that projects egocentric video segments and textual descriptions into a shared embedding space, incorporating the soft-label relevance matrices as supervision signal from training, with the architecture design validated by reproducing the published baseline scores on the validation split prior to any fine-tuning \textbf{(Weeks 3--4)}.

\item To implement and train the proposed model to estimate semantic similarity between video--text pairs and generate retrieval rankings for both search directions, using the SMS-Loss~\cite{wang2024symmetricmultisimilaritylossepickitchens100} and the AVION ViT-L/14~\cite{zhao2023avion} encoder as the starting point for experimentation, targeting a first evaluated checkpoint on the validation set \textbf{(Weeks 5--7)}.

\item To evaluate the model's performance using \textit{Normalized Discounted Cumulative Gain} (nDCG) and \textit{Mean Average Precision} (mAP) on the EK-100 MIR validation set, targeting results that match or exceed the official JPoSE + MI-MM challenge baseline~\cite{wray2019jpose,miech2020milnce} \textbf{(Weeks 7--8)}.

\end{enumerate}

\noindent\textit{Note: these objectives are to be reviewed and revised as needed once the first experimental results are available, as early findings may require adjusting scope, architecture choices, or evaluation targets. Weeks 8--9 are reserved for final analysis and written report.}